\section*{Resultados}

Nesta secção serão apresentados os resultados obtidos. Estes resultados devem conter (pelo menos) as seguintes subsecções:

\subsection*{Estudo de convergência do método dos painéis}

\begin{itemize}[itemsep=3pt]

\item Para um ângulo de ataque de \(5^\circ\), deve realizar-se uma análise que indique como varia o coeficiente de sustentação, \(c_l\), previsto pelo método dos painéis em função do número de painéis \(N\) utilizado. O resultado desta análise deve ser uma estimativa do número mínimo de painéis necessário para considerar que o parâmetro \(N\) já não afeta a solução.

\item Este estudo deve ser acompanhado por, pelo menos, um gráfico que mostre a evolução do coeficiente de sustentação em função do número de painéis.

\item Deve incluir-se uma imagem da painelização realizada após a escolha do número de painéis \(N\)

\end{itemize}

\subsection*{Estudo de convergência da simulação CFD}

\begin{itemize}[itemsep=3pt]

\item Deve realizar-se um estudo de convergência de malha e de independência da solução. Deve justificar a escolha da malha que considerou “ideal” para a realização da simulação. 

\end{itemize}

\subsection*{Comparação de métodos}

\begin{itemize}[itemsep=7pt]

\item Devem comparar-se os coeficientes de sustentação obtidos em função do ângulo de ataque utilizando os diferentes métodos. Para tal, deve fornecer-se, no mínimo, uma figura onde se sobreponham os coeficientes de sustentação previstos por cada um dos procedimentos.

\item Para o método dos perfis delgados, o ângulo de ataque deve variar entre \(-20^\circ\) e \(20^\circ\) com incrementos de \(0{,}5^\circ\), de modo a obter-se uma curva "contínua".

\item Para o método dos painéis, o ângulo deve variar entre \(-20^\circ\) e \(20^\circ\) com intervalos de \(5^\circ\).

\item Para a simulação CFD, devem analisar-se, no mínimo, os ângulos \(0^\circ\), \(5^\circ\), \(10^\circ\), \(15^\circ\) e \(20^\circ\).

\item Deve analisar-se a diferença entre os diferentes métodos com base nos resultados de \(c_l\). A análise dos vários procedimentos deve também ser feita em termos do tempo de cálculo e dos recursos necessários para cada um.

\item Deve identificar-se o ângulo de sustentação nula e o declive da curva de sustentação na zona linear previstos pelos três métodos.

% No início da terceira página, aproxima-se o espaçamento dos itens ao DOCX.
\setlength{\itemsep}{4pt}
\item Deve concluir-se quais as vantagens e desvantagens apresentadas por cada método e, com base nisso, explicar em que gama de ângulos de ataque ou para que aplicação se deveria empregar cada método.

\item Para os ângulos de \(0^\circ\), \(5^\circ\) e \(20^\circ\), devem apresentar-se as linhas de corrente previstas pelo método dos painéis e as linhas de corrente previstas pela simulação CFD.

\item Para os casos de \(0^\circ\), \(5^\circ\) e \(20^\circ\), deve apresentar-se uma comparação dos resultados previstos pelo método dos painéis e pela análise CFD, indicando e explicando as semelhanças e as diferenças. Para os diferentes ângulos, devem também identificar-se a posição do ponto de estagnação e o coeficiente de pressão mínimo previsto quer pelo método dos painéis, quer pela simulação CFD.

\item Com base nos resultados apresentados, deve indicar-se a gama de ângulos positivos na qual o perfil aerodinâmico apresenta um comportamento linear.

\end{itemize}
